\begin{trapbox}
	\begin{description}[style=nextline, leftmargin=2.8em, labelsep=0.5em, itemsep=0.35em]
		\item[\textbf{\textcolor{CTrap}{易错点一（忽略对称）}}] 直接连接$AB$求交点，认为$PA+PB$最小值为$AB$。
		\item[\textbf{\textcolor{CTrap}{正确理解（同侧需对称）：}}] 当$A,B$在$l$同侧时，必须通过作对称转化为异侧，最小值为$A'B$。
		\item[\textbf{\textcolor{CTrap}{错因分析（概念混淆）：}}] 混淆了同侧与异侧的处理方法，直接套用异侧模型。
	\end{description}

	\medskip
	\begin{description}[style=nextline, leftmargin=2.8em, labelsep=0.5em, itemsep=0.35em]
		\item[\textbf{\textcolor{CTrap}{易错点二（对称点错误）}}] 作对称点时，对称轴找错或计算错误，导致$A'$位置错误。
		\item[\textbf{\textcolor{CTrap}{正确理解（垂直平分）：}}] 对称点满足$l$垂直平分$AA'$，需利用垂直和中点条件求解。
		\item[\textbf{\textcolor{CTrap}{错因分析（几何性质不清）：}}] 对轴对称的几何条件掌握不牢，画图不准确。
	\end{description}

	\medskip
	\begin{description}[style=nextline, leftmargin=2.8em, labelsep=0.5em, itemsep=0.35em]
		\item[\textbf{\textcolor{CTrap}{易错点三（取等条件遗漏）}}] 认为最小值$A'B$在任何$P$都成立，忽略$P$需是$A'B$与$l$的交点。
		\item[\textbf{\textcolor{CTrap}{正确理解（三点共线）：}}] 仅当$A'$、$P$、$B$三点共线时取等，$P$是交点。
		\item[\textbf{\textcolor{CTrap}{错因分析（忽视放缩条件）：}}] 在应用不等式或几何公理时，忽视等号成立条件。
	\end{description}
\end{trapbox}
